\section{Список отозванных сертификатов}\label{FMT.CRL}

СОС издает УЦ, ранее выпустивший отзываемые сертификаты. Формат СОС описывается 
типом \code{CertificateList}, который определен в СТБ~34.101.19. 

Идентификатор алгоритмов ЭЦП, указываемый в компоненте~\code{signatureAlgorithm}
основного контейнера~\code{CertificateList} и дублируемый в
компоненте~\code{signature} вложенного контейнера~\code{TBSCertList}, должен
быть выбран из перечня, заданного в~\ref{CRYPTO.Sign}. Этот идентификатор должен
соответствовать открытому ключу издателя СОС (см.~\ref{CRYPTO.Keypair}).

Вложенный контейнер~\code{TBSCertList} заполняется по правилам СТБ~34.101.19 со
следующими уточнениями.

\begin{enumerate}
\item
Идентификационные данные издателя СОС, которые указываются в 
компоненте~\code{issuer}, должны повторять данные в одноименном 
компоненте его сертификата. 

\item
В каждую запись об отозванном сертификате (компонент~\code{revokedCertificates})
должно включаться расширение~\code{reasonCode} с кодом причины отзыва
сертификата.
%
В запись может включаться расширение~\code{invalidityDate}, которое описывает
момент наступления события, повлекшего отзыв.

\item
В расширении~\code{reasonCode} должны использоваться коды из следующего списка:
\begin{itemize}
\item
\code{unspecified}~--- неопределенная причина;
\item
\code{keyCompromise}~--- компрометация личного ключа субъекта (кроме ЦАС); 
\item
\code{cACompromise}~--- компрометация личного ключа УЦ;
\item
\code{affiliationChanged}~--- смена идентификационных данных субъекта;
\item
\code{superseded}~--- смена сертификата (например, с помощью~\code{Reenroll}, 
см.~\ref{PROCESSES.Reenroll});
\item
\code{cessationOfOperation}~--- прекращение деятельности субъекта;
\item
\code{privilegeWithdrawn}~--- отмена владения сертификатом (со стороны 
издателя);
\item
\code{aACompromise}~--- компрометация личного ключа ЦАС.
\end{itemize}

% info: https://security.stackexchange.com/questions/17432

\item
В список расширений СОС (компонент~\code{crlExtensions})
должны быть включены расширения~\code{AuthorityKeyIdentifier} 
и~\code{CRLNumber}. Первое расширение формируется по правилам,
заданным в~\ref{FMT.Ext.SKID}. Номер текущего СОС во втором расширении
должен быть неотрицательным целым числом, DER-код которого укладывается в 
20 октетов. Номера последовательных СОС должны монотонно возрастать.
\end{enumerate}

Список отозванных атрибутных сертификатов формируется по аналогичным правилам. 
Отличие в том, что список издает не УЦ, а ЦАС. Соответственно в компоненте 
\code{issuer} списка указываются идентификационные данные ЦАС.
